\documentclass[12pt]{article}
\everymath{\displaystyle}
\usepackage{unicode-math}
\usepackage{pgfplots}
\pgfplotsset{compat=1.18}
\usepackage{amsmath}



\usepackage[margin=1in]{geometry}
\usepackage{tcolorbox}
\usepackage{graphicx}
\usepackage{tikz}
\usepackage{amssymb}
\usepackage{enumitem}
\usepackage{titlesec}
\usepackage{fancyhdr}
\usepackage{pifont}
\usepackage{xcolor}
\newcounter{questionnum}
\setcounter{questionnum}{0}
\usepackage{eso-pic} % for page background

% --------- 主题颜色设置 ---------
\definecolor{novelblue}{RGB}{70,60,150}
\definecolor{headerbg}{RGB}{240,240,255}

\pagestyle{fancy}
\fancyhf{}
\renewcommand{\headrulewidth}{0.4pt}
\renewcommand{\headrule}{\hbox to\headwidth{\color{novelblue}\leaders\hrule height 0.4pt\hfill}}

% --------- 页眉背景色条 ---------
\newcommand{\headbackground}{
  \AddToShipoutPictureBG*{
    \begin{tikzpicture}[remember picture, overlay]
      \fill[blue!5!purple!10] (current page.north west) rectangle ([yshift=-60pt]current page.north east);
    \end{tikzpicture}
  }
}

% --------- 页眉设置 ---------
\pagestyle{fancy}
\fancyhf{}
\renewcommand{\headrulewidth}{0pt}
\fancyhead[L]{\color{novelblue}\textbf{\textbf{Novel Prep} }}
\fancyhead[C]{\color{novelblue}\textbf {SAT Preparation}}
\fancyhead[R]{\color{novelblue}\textbf{Jack Zeng}}


\newtcolorbox{SATCard}[1][]{%
  enhanced,
  colback=white,
  colframe=novelblue,
  coltitle=white,
  boxrule=0.6pt,
  arc=4pt,
  boxsep=5pt,
  left=8pt, right=8pt, top=10pt, bottom=10pt,
  sharp corners=south,
  drop shadow south east,
  #1
}

% --------- 顶部 Mark for Review 栏 ---------
\newcommand{\markforreview}{
  \stepcounter{questionnum} % 自动递增题号
  \begin{tikzpicture}[scale=1.1]
    % 黑色题号
    \fill[black] (0,0) rectangle (0.8,0.6);
    \node[white, font=\bfseries\large] at (0.4,0.3) {\thequestionnum};

    % 灰底主栏
    \fill[gray!10] (0.8,0) rectangle (14.5,0.6);

    % Mark for Review
    \node[anchor=west, font=\small] at (1.2,0.3) {\ding{72} \hspace{2pt}Mark for Review};

    % ABC 按钮区域
    \draw[thick, rounded corners=3pt] (13.5,0.12) rectangle (14.5,0.48);
    \node[font=\bfseries\normalsize] at (14.0,0.3) {ABC};
  \end{tikzpicture}


  \newcommand{\choiceboxcircled}[2]{%
  \begin{tcolorbox}[
    colback=white,
    colframe=black,
    boxrule=0.5pt,
    rounded corners=6pt,
    left=6pt, right=6pt, top=6pt, bottom=6pt,
    width=\linewidth
  ]
    \textbf{\large\textcircled{\textbf{#1}}} \ #2
  \end{tcolorbox}
}
}

% --------- 选项框样式 ---------
\newcommand{\choicebox}[2]{%
  \begin{tcolorbox}[
    colback=white, 
    colframe=novelblue,
    boxrule=0.5pt, 
    rounded corners=4pt, 
    left=6pt, right=6pt, top=6pt, bottom=6pt,
    width=\linewidth
  ]
  \textbf{#1.} \ #2
  \end{tcolorbox}
}


\begin{document}


\section*{SAT Math Hard Question 5}



\noindent\markforreview
\vspace{1em}

The given expression can be rewritten as \(\frac{a}{4}x^2 + \frac{b}{4}x + \frac{c}{4}\), where \(a, b,\) and \(c\) are constants. What is the value of \(a + b + c\)?  
Given the expression  
\[
-16(5x - 3)^2 + 4(5x - 2)^2
\]

\vspace{0.75em}

\choicebox{A}{-112}
\choicebox{B}{-56}
\choicebox{C}{-28}
\choicebox{D}{-7}

\vspace{2em}

\noindent\markforreview
\vspace{1em}

In the given function, \(b\) is a positive constant. For every increase in the value of \(x\) by 1, the value of \(f(x)\) increases by \(c\%\), where \(0 < c < 100\).  
Which expression gives the value of \(c\) in terms of \(b\)?

\[
f(x) = 66(b)^x
\]

\vspace{0.75em}

\choicebox{A}{1 + $\frac{b}{100}$}
\choicebox{B}{b + 100}
\choicebox{C}{100(b + 1)}
\choicebox{D}{100(b - 1)}

\vspace{2em}
\newpage
\noindent\markforreview
\vspace{1em}

If \(\sin(x^\circ) = a\), which of the following must be true for all values of \(x\)?

\vspace{0.75em}

\choicebox{A}{\cos(x^\circ) = a}
\choicebox{B}{\sin(90^\circ - x^\circ) = a}
\choicebox{C}{\sin(x^2)^\circ = a^2}
\choicebox{D}{\cos(90^\circ - x^\circ) = a}

\vspace{2em}

\noindent\markforreview
\vspace{1em}

The function \(g\) is a quadratic function in the \(xy\)-plane. The graph of \(y = g(x)\) has a vertex at \((-1, -4)\) and passes through the points \((-2, -43)\) and \((1, -160)\).  
What is the value of \(g(0) - g(2)\)?

\vspace{0.75em}

\choicebox{A}{312}
\choicebox{B}{355}
\choicebox{C}{-117}
\choicebox{D}{-203}

\vspace{2em}
\newpage
\noindent\markforreview
\vspace{1em}

The mass of object A is 324\% of the mass of object B and the mass of object A is 7.2\% of the mass of object C.  
If the mass of object C is \(p\%\) of the mass of object B, what is the value of \(p\%\)?

\vspace{0.75em}

\choicebox{A}{4.5\%}
\choicebox{B}{4.5}
\choicebox{C}{45}
\choicebox{D}{450}

\vspace{2em}

\noindent\markforreview
\vspace{1em}

In the given equation, \(c\) is a constant. If the equation has exactly one solution, what is the value of \(c\)?

\[
\frac{1}{cx} = \frac{x}{76} + \frac{1}{c}
\]

\vspace{0.75em}


\newpage

\noindent\markforreview
\vspace{1em}

A parallel electric circuit contains four resistors with resistances of $p$ ohms, $q$ ohms, $s$ ohms, and 13 ohms.  
The given equation relates the resistance $R$, in ohms, of the circuit, to the resistances of the four individual resistors it contains.  
Which equation correctly expresses $R$ in terms of $p$, $q$, and $s$?

\[
\frac{1}{R} = \frac{1}{p} + \frac{1}{q} + \frac{1}{s} + \frac{1}{13}
\]

\vspace{0.75em}

\choicebox{A}{$R = \frac{13pqs}{p+q+s+13}$}
\choicebox{B}{$R = \frac{13pqs}{13qs + 13ps + 13pq + pqs}$}
\choicebox{C}{$R = \frac{13qs + 13ps + 13pq + pqs}{13pqs}$}
\choicebox{D}{$R = p + q + s + 13$}

\vspace{2em}

\noindent\markforreview
\vspace{1em}

In the given equation, $b$ is an integer. The equation has no real solutions.  
What is the least possible value of $b$?

\[
-x^2 + bx - 169 = 0
\]

\vspace{0.75em}

\choicebox{A}{26}
\choicebox{B}{25}
\choicebox{C}{-26}
\choicebox{D}{-25}

\vspace{2em}

\noindent\markforreview
\vspace{1em}

The exponential function $f$ is defined by $f(x) = ab^x$, where $a$ and $b$ are positive constants.  
If $f(n-1) = f(n) + \frac{82}{100}f(n-1)$, where $n$ is a constant, what is the value of $b$?

\vspace{0.75em}

\choicebox{A}{1}
\choicebox{B}{0.82}
\choicebox{C}{0.18}
\choicebox{D}{10}

\vspace{2em}

\noindent\markforreview
\vspace{1em}

A rectangle is inscribed in a circle such that the length of the diagonal of the rectangle is twice the length of its shortest side.  
The circumference of the circle is $114\pi$ units. What is the area, in square units, of the rectangle?

\vspace{0.75em}

\choicebox{A}{$57\sqrt{2}$}
\choicebox{B}{$57\sqrt{3}$}
\choicebox{C}{$3,\!249\sqrt{2}$}
\choicebox{D}{$3,\!249\sqrt{3}$}

\vspace{2em}
\newpage
\noindent\markforreview
\vspace{1em}

In triangle $PQR$, the measure of angle $P$ is $(3x + 5)^\circ$, the measure of angle $Q$ is $(2x + 9)^\circ$, and the measure of angle $R$ is $(4y + 5)^\circ$.  
If side $QR$ is extended through point $R$ to point $S$, and the measure of angle $PRS$ is $(x + y)^\circ$, what is the value of $x + y$?

\vspace{0.75em}

\choicebox{A}{39}
\choicebox{B}{29}
\choicebox{C}{141}
\choicebox{D}{146}

\vspace{2em}

\noindent\markforreview
\vspace{1em}

The manager of a gym selected a sample of 139 members at random to estimate the percentage of the gym’s members that would continue to pay for a membership if the price increased.  
From the survey, the manager estimates that 82\% of the gym’s members would continue to pay for a membership if the price increased, with an associated margin of error of 6.39\%.  
If the survey is repeated with a random sample of 278 members and the results are calculated in the same way, which of the following will be the most likely effect of using the larger random sample compared to the smaller random sample?

\vspace{0.75em}

\choicebox{A}{The margin of error will be lower.}
\choicebox{B}{The margin of error will be higher.}
\choicebox{C}{The estimate of the percentage of the gym’s members that would continue to pay for a membership if the price increased will be lower.}
\choicebox{D}{The estimate of the percentage of the gym’s members that would continue to pay for a membership if the price increased will be higher.}

\vspace{2em}
\newpage
\noindent\markforreview
\vspace{1em}

Given the expression $\sqrt[5]{119n \left(\sqrt[6]{119n}\right)^2}$,  
for what value of $x$ is the expression equivalent to $(119n)^{30x}$, where $n > 1$?

\vspace{0.75em}

\choicebox{A}{225}
\choicebox{B}{$\frac{8}{225}$}
\choicebox{C}{$\frac{4}{225}$}
\choicebox{D}{$\frac{2}{225}$}

\vspace{2em}

\noindent\markforreview
\vspace{1em}

The mass of object A is 483\% of the mass of object B, and the mass of object A is 0.069\% of the mass of object C.  
If the mass of object C is $p\%$ of the mass of object B, what is the value of $\frac{p}{1,\!000}$?

\vspace{0.75em}

\choicebox{A}{0.7\%}
\choicebox{B}{70}
\choicebox{C}{70\%}
\choicebox{D}{700}

\vspace{2em}
\newpage
\noindent\markforreview
\vspace{1em}

For triangle $PQR$ and triangle $STU$, angles $P$ and $S$ each measure $30^\circ$, $QR = 16$ and $TU = 12$.  
Which additional piece(s) of information is(are) sufficient to prove that triangle $PQR$ is similar to triangle $STU$?

\begin{itemize}
    \item[I] Q = $40^\circ$ and U =$110^\circ$ 
    \item[II] $\frac{PR}{SU} = \frac{4}{3}$
\end{itemize}
\vspace{0.75em}

\choicebox{A}{I is sufficient, but II is not.}
\choicebox{B}{II is sufficient, but I is not.}
\choicebox{C}{I is sufficient, and II is sufficient.}
\choicebox{D}{Neither I nor II is sufficient.}



\end{document}
